\section{The intensity derivative and the pivotal amplitude}
\label{sec:amplitude}

The preceding section supplies two estimates with different roles.
Corollary~\ref{tr:insertion} makes the effect of a distant fair insertion
integrable, even after division by the central pivotal probability.
Corollary~\ref{tr:wedge} makes the normalized effect of a nearby insertion
almost independent of the size of the rectangle. Together they give a
summable error for the logarithmic derivative of $f$.
Throughout this section the Poisson intensity is one, unless otherwise
specified, and $\lambda=5/4$.

\subsection{Differentiating the intensity}

Let $X_R(\eta)$ be the indicator of a black left--right crossing of
$RQ_0$ for a colored point configuration $\eta$, with the closed-cell
convention of Theorem~\ref{thm:main}. If $0^{\rm b}$ and $0^{\rm w}$
denote an inserted black and white point at the origin, respectively, put
\[
 H_R(\eta)
 =X_R(\eta\cup\{0^{\rm b}\})-X_R(\eta\cup\{0^{\rm w}\}).
\]
Thus $H_R\in\{0,1\}$, and $f(R)=\E H_R(\eta)$ for the unit-intensity
fairly colored process. For $y\ne0$, write $y^{\rm f}$ for an independently
fairly colored inserted point and define
\begin{equation}\label{cl:kernel}
 K_R(y)=\E\bigl[H_R(\eta\cup\{y^{\rm f}\})-H_R(\eta)\bigr].
\end{equation}
The expectation includes the extra color. Values at $y=0$ and at null
nongeneric insertion locations are immaterial to the integrals below.
In particular, $K_R$ is an insertion difference of a \emph{color}-pivotal
probability, not a probability of Poisson pivotality.

\begin{lemma}[Intensity differentiation]\label{cl:derivative}
The function $f$ is continuously differentiable on $(0,\infty)$. For
every $R>0$, $K_R$ is absolutely integrable and
\begin{equation}\label{cl:derivative-eq}
 \frac R2 f'(R)=\int_{\R^2}K_R(y)\,dy.
\end{equation}
\end{lemma}

\begin{proof}
Let $\eta_t$ have intensity $t$ and fair colors. Euclidean scaling gives
\begin{equation}\label{cl:scaling-intensity}
 G_R(t):=\E H_R(\eta_t)=f(\sqrt t\,R).
\end{equation}
We justify differentiation of the full-plane expectation before using
this identity. All the estimates in this paragraph are uniform for
$t$ in a compact subinterval of $(0,\infty)$.

For $n\ge n_0(R)$, let $H_{R,n}$ be the same test computed from the
points in $B(0,2^n)$, together with the central insertion. A locality
screen in a collar well inside $B(0,2^n)$ determines every cell meeting
$RQ_0$, including its intersections with the boundary. Consequently
$H_{R,n}=H_R$ outside an event of probability at most
$C_R\exp(-c_R4^n)$. The same screen gives
$H_{R,n+1}=H_{R,n}$. Finite-volume Poisson differentiation follows by
differentiating the absolutely convergent Poisson series; see also
\cite[Theorem~19.1]{LP}. Equivalently,
for a bounded function $F$ of the process in a bounded set $D$,
\begin{equation}\label{cl:poisson-score}
 \frac{d}{dt}\E F(\eta_t\cap D)
 =\E\left[F(\eta_t\cap D)
       \left(\frac{N_t(D)}t-|D|\right)\right].
\end{equation}
Here $N_t(D)$ is the number of Poisson points in $D$.
Apply this to $F=H_{R,n+1}-H_{R,n}$ on $D=B(0,2^{n+1})$.
The score in parentheses has second moment $|D|/t$. Cauchy--Schwarz
therefore bounds the derivative by
$C_R2^n\exp(-c_R4^n)$. The differences of the expectations and of
their derivatives are uniformly summable. Hence $G_R$ is $C^1$ and
its derivative is the limit of the finite-volume derivatives.

The add-one version of finite-volume differentiation is
\[
 \frac{d}{dt}\E H_{R,n}(\eta_t)
 =\int_{B(0,2^n)}
  \E\bigl[H_{R,n}(\eta_t\cup\{y^{\rm f}\})
                     -H_{R,n}(\eta_t)\bigr]\,dy.
\]
For fixed $y$ the integrand converges to its full-plane counterpart.
For $|y|\ge C R+C$, a screen surrounding $RQ_0$ and the origin, at
distance comparable to $|y|$ and strictly short of $y$, prevents the
insertion from changing any cell used by the test. Its failure
probability is at most $C_R e^{-c_R|y|^2}$. The screen can be taken
inside $B(0,2^n)$ whenever $y\in B(0,2^n)$ and $n$ is sufficiently
large; thus this is also a dominating bound for the finite-volume
integrands. On a fixed bounded set the bound $2$ suffices.
Dominated convergence proves absolute integrability and identifies
$G_R'(1)$ with the integral in \eqref{cl:derivative-eq}.
Equation~\eqref{cl:scaling-intensity} proves the asserted derivative
formula and, locally at each $R$, the $C^1$ assertion for $f$.
\end{proof}

\subsection{An integrable tail after normalization}

The exponential localization in Lemma~\ref{cl:derivative} is a
finite-scale statement. We need a bound uniform as $R$ grows, with
the small factor $f(R)$ retained. This is where the signed insertion
estimate is essential.

\begin{lemma}[Normalized insertion tail]\label{cl:tail}
There are constants $\xi>0$, $m_0<\infty$, and $C<\infty$ such that
for $m_0\le m\le R/100$,
\begin{equation}\label{cl:tail-eq}
 \int_{|y|>m}|K_R(y)|\,dy
 \le C f(R)\bigl(m^{-\xi}+R^{-\xi}\bigr).
\end{equation}
\end{lemma}

\begin{proof}
Use the exponent $0<\nu<1$ of Corollary~\ref{tr:insertion} and
the one-arm exponent $\beta>0$ from the critical inputs. All screens
and arm probes below have fixed-factor buffers. Independent probes
mean probes determined by disjoint sets of Poisson coordinates,
not unrestricted Voronoi events in disjoint planar regions.

We first explain how arm requirements can be multiplied by a signed
message estimate. The integrand in \eqref{cl:kernel} is a mixed
difference: one difference changes the central color, and the other
inserts the point at $y$. If this mixed difference is nonzero, changing
the central color changes a crossing for at least one setting of the
insertion, and changing the insertion changes a crossing for at least
one setting of the central color. A screened probe separated from
both changes reads the same base colors in all four configurations.
Thus the necessary-arm implications for a local change apply to
these common colors; see \cite[Section~2, ``Necessary arms for a
local change'']{V}. They apply even when the two crossing changes
occur in different configurations.

The conditional version is equally important. Every formal noncrossing
pairing can be drawn as disjoint wall arcs in a topological disk.
Coloring its complementary regions alternately realizes the prescribed
boundary runs and the pairing. Thus two formal cap states give two
planar fillings agreeing on every exterior probe; the deterministic
necessary-arm argument does not require Poisson geometry inside the
replaced disk. Fix a successful
cut about $y$, its seam, and all exterior variables. The exterior
test is now a function of the inside connection state. If a
necessary central or enclosing four-arm event does not occur,
the central black-minus-white test is zero for every such state.
If a necessary side or corner arm event does not occur, each
crossing test with the central color fixed is constant across
inside states. These are
the same planar necessary-arm implications, applied to two
possible cap states on the cut. The insertion message has zero
mass, so it annihilates the latter constants. Its conditional
answer is consequently bounded by its total variation norm times
the necessary probe indicators. Those indicators are measurable
outside the cut. Disjoint buffered supports then allow their
probabilities to be multiplied.

We use screened arm probes with failures allowed, whose probabilities
are still bounded by a constant times their intended arm bounds;
see the convention in Section~\ref{sec:inputs} and
\cite[Section~2, ``Buffered localization'']{V}.
Indeed, group failures by the largest failed radius $v$ and retain
the arm test starting a fixed factor beyond $v$, independently of
the failure. For a four-arm probe starting at $r$, the relevant
estimate is
\[
 e^{-cv^2}A(C_0v,R)
 \le C A(r,R)(v/r)^g e^{-cv^2}
 \qquad(C_0v<R),
\]
where $g$ is the polynomial exponent in \eqref{in:power}.
For a side or corner probe the corresponding endpoint losses
are bounded by $C(v/d)^2$ or $C(v/u)^\beta$, respectively.
If a failed screen crosses the junction between side and corner
probes, the lost factors are therefore only polynomial. If the
retained test has no remaining radial range, the same comparison follows from the
polynomial lower bound, or directly from the displayed power
factor for a side or corner probe. Summing over geometric $v$
preserves the indicated arm cost as long as the retained probes
and the failure event have disjoint buffered supports. A failure
large enough to reach the central probe requires a separate
argument below: that probe is then dropped, not treated as
independent. An unweighted exponentially small error alone would
not suffice when $R$ is arbitrarily large.

\paragraph{Insertions near the center.}
Put $t=|y|$ and suppose $m<t\le R/100$. First fix a small
constant $c_1>0$, and then choose $a>0$ sufficiently small relative
to $c_1$. Choose a cut about $y$ of radius $a t$,
a central four-arm probe from $b_*$ to $a t$,
and a four-arm probe surrounding both insertions from $a^{-1}t$
to $aR$.
Short or overlapping endpoint ranges are omitted at constant cost.
The buffers of these three constructions can be chosen disjoint.
On cut success, Corollary~\ref{tr:insertion} gives a bound
$Ct^{-2-\nu}$ for the insertion message. The central color
difference requires the first four-arm probe; transmission of a
change from the region containing both insertions to the crossing
requires the second. Thus the success contribution is at most
$Ct^{-2-\nu}A(b_*,t)A(t,R)$.

For failures, let $v$ be the largest failed radius about $y$, with
the initial radius comparable to $t$. If $v<c_1t$, the
central probe and an enclosing
probe retain disjoint supports after the fixed buffer adjustments.
The preceding failed-screen estimate bounds their total
contribution by $Ce^{-ct^2}A(b_*,t)A(t,R)$.
If $v\ge c_1t$, drop the central probe. Retain instead a screened
four-arm probe centered at the origin, from radius $C_1(v+t)$
to $aR$, enclosing both modifications and with its support
separated from the failed screen. The same necessary-arm
implication supplies this probe. When its radial range is nonempty,
quasi-multiplicativity and \eqref{in:power} give
\[
 e^{-cv^2}A(C_1(v+t),aR)
 \le C A(b_*,R)(v+t)^g e^{-cv^2}.
\]
Since $t\le c_1^{-1}v$, summation over these geometric $v$ is
at most $C A(b_*,R)e^{-c't^2}$. If the retained radial range
is empty, then $R\le C(v+t)\le C'v$; the absolute failure
bound and $A(b_*,R)\ge cR^{-g}$ give the same conclusion.
Combining success and both failure cases proves
\begin{align}
 |K_R(y)|
 &\le Ct^{-2-\nu}A(b_*,t)A(t,R)
             +C A(b_*,R)e^{-ct^2}\notag\\
 &\le C f(R)t^{-2-\nu}.
 \label{cl:bulk-kernel}
\end{align}
In the last line we used quasi-multiplicativity and
\eqref{in:comparison}. In particular,
\begin{equation}\label{cl:bulk-integral}
 \int_{m<|y|\le R/100}|K_R(y)|\,dy
 \le C f(R)\int_m^\infty t^{-1-\nu}\,dt
 \le C f(R)m^{-\nu}.
\end{equation}

\paragraph{Insertions near a side or corner.}
Consider $R/100<|y|\le10R$, including points outside the rectangle.
Fix a sufficiently large constant $d_*$, and then a sufficiently
small constant $\rho>0$. If $\mathcal C_R$ is the set of the four
corners of $RQ_0$, set
\[
 d(y)=\min\{\rho R,\max\{d_*,\dist(y,\partial(RQ_0))\}\},
 \qquad
 u(y)=\min\{\rho R,\max\{d_*,\dist(y,\mathcal C_R)\}\}.
\]
We take $R$ large enough that $d_*<\rho R$. Always
$d_*\le d(y)\le u(y)\le\rho R$; write $d=d(y)$ and $u=u(y)$.
The constants in the estimates may depend on $d_*$ and $\rho$.
There is a central four-arm probe up to radius $cR$, where $c$
is fixed and much smaller than $1/100$, disjoint from the other
regions used below. Its cost is at most $CA(b_*,R)$.

At an interior point $y$, use the insertion estimate on a cut of
radius a small fixed fraction of $d$. At the bottom cutoff a
constant bound suffices, and can be bounded by $Cd^{-2-\nu}$
because $d_*$ is fixed. At an exterior point with distance above
the bottom cutoff, successful screening at this scale leaves the
quad unchanged. Only screen failures contribute there, with a bound
exponential in $d^2$, which is no larger than $Cd^{-2-\nu}$.

If $u/d$ is larger than a sufficiently large constant, the closest
part of the boundary is a straight side. A crossing change then
requires three alternating free arms in a half-plane probe centered
at the projection onto that side, from radius $Cd$ to radius $cu$.
Its cost is at most $C(d/u)^2$. If the projection is at a corner,
the two distances are comparable and this probe is unnecessary.
When $u$ is uncapped and is a sufficiently small fraction of $R$,
retain a one-arm probe centered at the nearest corner from radius
$Cu$ to radius $cR$. Its cost is at most $C(u/R)^\beta$.
For $u$ capped above, this factor is a fixed positive constant,
so the corresponding probe can be omitted. These are precisely
the local-change arm implications for the unwired, closed-cell
crossing convention. They also cover a modified cell meeting a
side or a corner.

The insertion cut, the side probe, the corner probe, and the
central probe can be given disjoint buffered supports. To arrange
this, first fix the central radius $cR$, then choose the distance
cap $\rho R$ much smaller. A corner probe is used only in a
corner neighborhood, far from the center. Fixed multiples are
lost between consecutive probes, which changes only constants.
To include failures, choose a sufficiently small fixed $c_2>0$.
For a largest failed radius $v<c_2R$ about $y$, the central
probe remains disjoint from the failed-screen support. Retain
the side and corner probes only beyond that support. If a
junction has been passed, the center of the affected side or
corner stage lies within $O(v)$ of $y$, so a fixed-factor
enlargement provides the required separation. Omitting an entire
side stage loses at most $C(1+v/d)^2$, and shifting or omitting
a corner stage loses at most $C(1+v/u)^\beta$. Since the first
possible failed radius is a fixed fraction of $d$,
\[
 \sum_{v\ge c_3d}e^{-cv^2}(1+v/d)^2(1+v/u)^\beta
 \le Cd^{-2-\nu}.
\]
The retained central probe supplies $A(b_*,R)$ throughout this
case. For $v\ge c_2R$, discard all arm probes and sum the
absolute failure probabilities, obtaining $Ce^{-c'R^2}$.
This is bounded by the right side of \eqref{cl:boundary-kernel}:
$A(b_*,R)\ge cR^{-g}$ and
$d_*\le d\le u\le\rho R$ make the reciprocal of that right
side at most a fixed power of $R$. This case includes every
failure whose support reaches the central probe, because
$|y|\ge R/100$ and the central probe radius was chosen much
smaller than $R/100$.
Thus no overlapping supports have been treated as independent.
The conditional total-variation argument on screen success,
together with these failure bounds, proves uniformly that
\begin{equation}\label{cl:boundary-kernel}
 |K_R(y)|\le C A(b_*,R)\,
     d^{-2-\nu}\left(\frac d u\right)^2
                    \left(\frac u R\right)^\beta.
\end{equation}

For completeness, a dyadic class with distances comparable to
$d,u$ has area at most $Cud$. Below the upper cap, it lies in
corner disks of radius $Cu$ and boundary strips of width $Cd$.
If $u$ is capped, the total strip length is $O(R)=O(u)$;
if both distances are capped, the ambient region has area
$O(R^2)=O(ud)$. These assertions apply equally outside the
rectangle. The bottom classes have the same bounds with $d_*$
in place of a smaller distance. Summing
\eqref{cl:boundary-kernel} first over dyadic $d\le u$ gives
\begin{align*}
 \sum_{d\le u} Cud\,d^{-2-\nu}(d/u)^2(u/R)^\beta
 &=\frac C u\left(\frac uR\right)^\beta
                           \sum_{d\le u}d^{1-\nu}\\
 &\le C u^{-\nu}(u/R)^\beta,
\end{align*}
where $\nu<1$ is used. The remaining dyadic sum is bounded by
\[
 C R^{-\beta}\sum_{d_*\le u\le\rho R}u^{\beta-\nu}
 \le C R^{-\xi_0}
 \qquad\text{for every }0<\xi_0<\min\{\beta,\nu\}.
\]
The strict inequality accommodates the logarithm when $\beta=\nu$.
Together with \eqref{in:comparison}, this proves
\begin{equation}\label{cl:boundary-integral}
 \int_{R/100<|y|\le10R}|K_R(y)|\,dy
 \le C f(R)R^{-\xi_0}.
\end{equation}

\paragraph{Far exterior insertions.}
For $|y|>10R$ an agreement screen between the test rectangle and
$y$ gives $|K_R(y)|\le Ce^{-c|y|^2}$. Its integral is at most
$Ce^{-cR^2}$. The polynomial lower bound for $f$ following from
\eqref{in:comparison} absorbs this into $Cf(R)R^{-\xi_0}$.
Combining this with \eqref{cl:bulk-integral} and
\eqref{cl:boundary-integral}, and decreasing the exponent, proves
\eqref{cl:tail-eq}.
\end{proof}

\subsection{Comparing the normalized derivatives}

The tail estimate permits us to discard insertions far from the
origin. For the remaining insertions we use the same interior
messages in two different rectangles. Rank-one approximation then
controls a determinant, so that no unknown leading factor needs
to be identified.

\begin{lemma}[Two-size determinant]\label{cl:determinant}
Let $\nu'>0$ be the exponent in Corollary~\ref{tr:wedge}.
There are constants $C,c>0$ such that, for sufficiently large $m$,
$100m\le R\le R'\le2R$, and almost every $0<|y|\le m$,
\begin{equation}\label{cl:determinant-eq}
 \bigl|K_R(y)f(R')-K_{R'}(y)f(R)\bigr|
 \le C(m/R)^{2\lambda+\nu'}+Ce^{-cm^2}.
\end{equation}
The constants are independent of the distance from $y$ to the origin.
\end{lemma}

\begin{proof}
Take concentric cuts about the origin, starting at radius $10m$
and ending at a radius comparable with a small fixed fraction of
$R$, with consecutive radii in ratio $L$. Round the final radius
down to the last such scale. Its ratio to $R$ is bounded below
by a positive constant depending only on the already fixed $L$.
All cut collars lie strictly inside both rectangles and outside
both inserted points. If necessary, enlarge the fixed lower
bound on $R/m$; the remaining bounded ratios follow by increasing
$C$, since all four scalar quantities in the determinant are bounded.

First restrict every test to success of every cut screen, with
the screens evaluated on the common base process. The probability
discarded is bounded by
\[
 C\sum_{j\ge0}\exp\bigl(-cm^2L^{2j}\bigr)
 \le Ce^{-c'm^2}.
\]
At the first cut let $a$ be the signed connection-state message
obtained by changing the central color. Let $b_y$ be the
difference of this message with and without the extra fair point
at $y$. Both have zero mass conditional on each seam, and
their total variation norms are bounded by $2$ and $4$,
respectively. They depend on neither $R$ nor $R'$: on screen
success their construction uses only the common interior
Poisson variables and the first seam. In particular this
description does not require disjoint neighborhoods of $0$ and $y$.

Let $T$ be the composed transfer through the intervening cuts.
The exterior crossing tests in $RQ_0$ and $R'Q_0$ define bounded
linear functionals $\ell_R$ and $\ell_{R'}$ at the final cut.
The four restricted expectations are
\[
 \ell_R(Ta),\quad \ell_{R'}(Ta),\quad
 \ell_R(Tb_y),\quad \ell_{R'}(Tb_y).
\]
Their determinant is the evaluation of $T^{\otimes2}$ on
$b_y\otimes a-a\otimes b_y$ by
$\ell_R\otimes\ell_{R'}$. Corollary~\ref{tr:wedge} therefore
bounds it by $C(m/R)^{2\lambda+\nu'}$.
Restoring the discarded events changes the determinant by at
most $Ce^{-cm^2}$, because each scalar factor is bounded.
This proves \eqref{cl:determinant-eq}.
\end{proof}

\begin{proposition}[Summable logarithmic error]\label{cl:log-derivative}
There are constants $\alpha>0$ and $C<\infty$ such that
\begin{equation}\label{cl:log-derivative-eq}
 \left|\frac{Rf'(R)}{f(R)}+\lambda\right|
 \le CR^{-\alpha}
\end{equation}
for all sufficiently large $R$.
\end{proposition}

\begin{proof}
Set $g(R)=Rf'(R)/f(R)$, which is defined and continuous for
large $R$ by Lemma~\ref{cl:derivative} and
\eqref{in:comparison}. We first prove a uniform estimate on
each interval $[R,2R]$. Fix
\[
 \zeta=\frac{\nu'}{4(2+2\lambda+\nu')},\qquad
 \eta=\frac{\nu'}8,\qquad m=R^\zeta.
\]
The exponent estimate \eqref{in:exponent-eq} implies
$f(s)\ge s^{-\lambda-\eta}$ for all sufficiently large $s$.
For $R\le R'\le2R$, divide \eqref{cl:determinant-eq} by
$f(R)f(R')$ and integrate over $|y|\le m$. The result is
\begin{align}
 \left|\int_{|y|\le m}\frac{K_R(y)}{f(R)}\,dy
       -\int_{|y|\le m}\frac{K_{R'}(y)}{f(R')}\,dy\right|
 &\le C R^{2\lambda+2\eta}m^2
          \bigl[(m/R)^{2\lambda+\nu'}+e^{-cm^2}\bigr]\notag\\
 &\le C R^{-\nu'/2}.
 \label{cl:normalized-near}
\end{align}
Indeed the power in the first term is exactly
\[
 -\nu'+2\eta+\zeta(2+2\lambda+\nu')=-\nu'/2,
\]
and the exponential term is smaller. For large $R$, the chosen
$m$ satisfies all the scale restrictions of
Lemmas~\ref{cl:tail} and~\ref{cl:determinant}.
Lemma~\ref{cl:tail} bounds each complementary normalized integral
by $C R^{-\zeta\xi}$. Since
$g(R)=2\int K_R(y)/f(R)\,dy$ by
\eqref{cl:derivative-eq}, we obtain
\begin{equation}\label{cl:dyadic-oscillation}
 \sup_{R\le R'\le2R}|g(R')-g(R)|\le CR^{-\alpha},
 \qquad
 \alpha=\min\{\nu'/2,\zeta\xi\}>0.
\end{equation}

The values $g(2^nR_0)$ form a Cauchy sequence, because the
successive errors in \eqref{cl:dyadic-oscillation} are summable.
The same estimate between dyadic points shows that $g(R)$ has
a finite limit $\gamma$ as $R\to\infty$, and summing along
$R,2R,4R,\ldots$ gives
\[
 |g(R)-\gamma|\le C\sum_{j\ge0}(2^jR)^{-\alpha}
 \le C'R^{-\alpha}.
\]
It remains to identify this limit. Integration yields
\[
 \log f(R)-\log f(R_0)=\int_{R_0}^R g(s)\,\frac{ds}s.
\]
Dividing by $\log R$ shows that
$\log f(R)/\log R\to\gamma$. By \eqref{in:exponent-eq} the
same limit is $-\lambda$. Thus $\gamma=-\lambda$, proving
\eqref{cl:log-derivative-eq}.
\end{proof}

\subsection{Completion of the normalization}

\begin{proof}[Proof of Theorem~\ref{thm:main}]
By Proposition~\ref{cl:log-derivative},
\[
 \frac{d}{dR}\log\bigl(R^\lambda f(R)\bigr)=O(R^{-1-\alpha}).
\]
The right side is integrable at infinity. Since $f(R)>0$ for
large $R$, the logarithm has a finite real limit. Consequently
\begin{equation}\label{cl:amplitude}
 R^{5/4}f(R)\longrightarrow a\in(0,\infty).
\end{equation}
Positivity and finiteness follow from convergence of the
\emph{logarithm}, rather than merely boundedness of the normalized
probabilities. In fact the argument also gives
$R^{5/4}f(R)=a(1+O(R^{-\alpha}))$.

Recall that $S_\eps=f(1/\eps)$,
$m_\eps=\eps^{-2}$, and
$q_\eps=(\E N_\eps(Q_0))^{-1}$.
The calibration limits \eqref{in:calibration-eq} are
\[
 m_\eps q_\eps S_\eps\longrightarrow c\in(0,\infty),
 \qquad
 q_\eps\E N_\eps\longrightarrow B\in(0,\infty).
\]
Equation~\eqref{cl:amplitude} gives
$S_\eps\sim a\eps^{5/4}$, hence
$q_\eps\sim(c/a)\eps^{3/4}$. It follows that
\[
 \eps^{3/4}\E N_\eps\longrightarrow \frac{Ba}{c}
 =:c_V\in(0,\infty).
\]
Here $N_\eps$ counts color-pivotal sites for the original unwired
unit-square event. The Poisson insertion was used only to
differentiate $f$; it has not changed the pivotal count in this
last formula. Taking reciprocals proves the equivalent assertion
for the near-critical scale.
\end{proof}
